% Nouveau EGA V (cote 25): archive cote 25, site page 61 (PDF page 62).
% typescript. Draft compared: none.
% AI-transcribed by Claude Opus 5.5 (claude-opus-5-5, adaptive thinking, effort high), Message Batches API, batch ega25b (msgbatch_01A7n5YuXubUPNK1zsepi8RG), 2026-11-13; checked against both supplied scan parts; see CHECK_DECISIONS.tsv.
\scanpage{25/61}
\mspage{56}

L'énoncé 10.2. est une conséquence à peu près immédiate du suivant,
\placedmargin{|}qui est tout à fait indépendant d'histoires de sections hyperplanes, et
serait sans doute mieux à sa place dans  un nº en rabiot aux
``immersions régulières''.

\placedmargin{|}\placedmargin{\unsure{$\backslash\,/$}}\uline{Proposition} 10.3. Soient X un préschéma,  $\underline{G}$ un Module quasi-cohérent
sur X, et $u\colon \underline{G}\hand{\to}\underline{O}_X$ un homomorphisme, $\underline{J} = u(G)$, $T=V(\underline{J})$. Soit
$\overset{\hand{\sim}}{X}$ déduit de X par éclatement de T. Considérons d'autre part
$P = P(\underline{G})$, l'homomorphisme canonique $\underline{G}_P\hand{\to}\underline{O}_P(1)$ et son noyau
$\underline{H}$ ,(de sorte que nous avons la suite exacte $0\hand{\to}\underline{H}\hand{\unsure{\relbar|\to}}\underline{G}_P\hand{\to}\underline{O}_P(1)\hand{\to}0$),
l'homomorphisme $u_P : \underline{G}_P \hand{\to} \underline{O}_P$ , et l'Idéa\unsure{l|m} qu coh $\underline{K} = u_P(\underline{H})\hand{\subset}\underline{O}_P\placedmargin{\unsure{\smile}}$ .
Alors $\overset{\hand{\sim}}{X}$ est canoniquement isomorphe à un sous-préschéma fermé de
$V(\underline{K})$ . Lorsque G est localement libre et u est ``régulier'', alors
l'isomorphisme précédent est un isomorphisme de $\overset{\hand{\sim}}{X}$ avec $V(\underline{K})$ lui-même,
de plus dans ce cas $\underline{H}$ est\placedmargin{\unsure{o|\illegible{1}}} localement libre sur P et  $\underline{H}\hand{\to}\underline{O}_P$
(dont le préschéma d es zéros e st $\overset{\hand{\sim}}{X}$ ) est également régulier.

Le premier énoncé est à peu près trivial. Le deuxiè\unsure{m|p}e est un
exercice qui ne peut offrir de difficulté (je ne l'ai pas fait en
détail, pensant que tu t'en tireras aussi bien que moi).

Lorsque dans 10.2.  $S = \mathrm{Spec}\,k$ et X est de dimension $\hand{\leqslant}1$,
alors l'hypothèse de régularité faite équivaut à $T = \hand{\unsure{\emptyset|\phi}}$ , de sorte
que $Y_{\underline{L}}\hand{\to} X$ est un \uline{isomorphisme} . On trouve ainsi\hand{;} en conjuguant
avec 9.2. :

\uline{Corollaire} 10.4. (à 10.2.) Soit X une courbe lisse géom connexe dans
 l'espace projectif P sur un corps infini k, et soit $n\hand{\geqslant}2$.
Alors il existe un pinceau linéaire de n-formes sur P, définissant un
morphisme $X\hand{\to}P^1$ ayant les propriétés suivantes: ce morphisme est

\mspage{61}
